Legacy systémy v jazyku Java predstavujú pre vývojárov dlhodobý problém: chýba im dokumentácia aj automatizované testy, a každá zmena je preto riziková.
Zavedenie jednotkových testov by toto riziko znížilo, no ich spätné vytváranie je časovo náročné a často naráža na zlú testovateľnosť kódu,
najmä na silné závislosti medzi komponentmi.
Cieľom diplomovej práce je preskúmať možnosti využitia veľkých jazykových modelov (LLM) na automatizované generovanie jednotkových testov pre takéto systémy. 
V práci navrhujeme metodiku na generovanie a vyhodnocovanie testov a prakticky overujeme schopnosti viacerých LLM modelov.
Modely porovnávame empiricky z hľadiska úspešnosti generovania, pokrytia kódu a spoľahlivosti vygenerovaných testov.
Súčasťou práce je aj analýza vplyvu štrukturálnych vlastností zdrojového kódu, ako sú cyklomatická a kognitívna komplexita či miera závislostí, na jeho testovateľnosť a úspešnosť generovania testov.
Zároveň skúmame, či dokáže AI kód upraviť tak, aby bol lepšie testovateľný. 
% Výsledky ukázali, že [doplniť hlavné zistenia, napr. ktorý model dosiahol najlepšie pokrytie a ako komplexita ovplyvňuje úspešnosť].
% Práca poskytuje [doplniť prínos, napr. odporúčania pre praktické nasadenie LLM pri testovaní legacy kódu].